\section{接口设计}

\subsection{接口设计目标}
本系统接口设计的目标，是在当前项目已有的用户端、管理端、后端业务服务、异步任务调度和实时通知基础上，形成统一、清晰、可扩展的接口体系。接口设计既要服务当前已实现的图像上传、图像检测、人工审核、通知与管理功能，也要为论文检测、Review 检测、综合鉴伪报告和模型扩展预留统一接入方式。与之配套，用户端需配置清晰的页面路由与导航结构（见本章第~\ref{sec:user-client-routing}~节），使 REST/WebSocket 能力与界面入口一一可追溯。

从系统边界上看，本章接口设计主要覆盖三类接口：
\begin{itemize}[leftmargin=2em]
    \item 面向前端用户端和管理端的 HTTP REST 接口；
    \item 面向前端的 WebSocket 实时推送接口；
    \item 面向系统内部的异步任务调度接口、AI 推理调用接口和文件访问接口。
\end{itemize}

\subsection{接口设计原则}
\begin{itemize}[leftmargin=2em]
    \item \textbf{统一入口原则}：系统外部接口统一经由后端业务服务对外暴露，前端不直接访问数据库或 AI 推理服务；
    \item \textbf{REST 风格优先}：外部 HTTP 接口优先采用资源路径加标准方法的方式组织，便于前后端协作与后续接口治理；
    \item \textbf{同步与异步分离}：上传、查询、审批等短时操作采用同步接口；检测任务、报告生成等长时操作通过异步任务接口和状态查询接口解耦；
    \item \textbf{统一任务、差异结果}：接口层对不同检测对象尽量采用统一任务入口、统一状态查询接口，并通过差异化结果接口返回图像、论文和 Review 特有数据；
    \item \textbf{权限前置}：用户认证、角色检查、组织边界和资源归属校验应在接口层或接口进入业务层前完成；
    \item \textbf{可演进性}：接口命名和路径设计应兼顾当前实现与后续统一资源抽象，不将接口永久绑定到单一图像对象模型上。
\end{itemize}

\subsection{接口总体分层}
从调用关系看，本系统接口可划分为四层：
\begin{longtable}{|p{3cm}|p{10.8cm}|}
\hline
层次 & 说明 \\ \hline
前端访问层 & 面向用户端与管理端提供 HTTP API 和 WebSocket 接口，负责请求接入、参数传递和结果展示。 \\ \hline
业务接口层 & 由 Django + DRF 暴露具体接口，负责认证鉴权、参数校验、业务编排和响应封装。 \\ \hline
系统协作层 & 负责后端与 Celery 任务、WebSocket 通道、文件系统和报告生成逻辑之间的协作。 \\ \hline
能力适配层 & 负责后端与具体 AI 推理服务、图像检测算法、论文检测模型和 Review 检测模型之间的调用适配。 \\ \hline
\end{longtable}

\subsection{外部接口分类}
当前项目外部接口可分为以下几类：
\begin{itemize}[leftmargin=2em]
    \item 用户认证与个人信息接口；
    \item 组织与邀请码接口；
    \item 资源上传、查询与资源管理接口；
    \item 检测任务提交、状态查询、结果查询和报告下载接口；
    \item 人工审核申请、审核执行与审核查询接口；
    \item 通知与消息接口；
    \item 管理端用户管理、资源治理、仪表盘统计和日志管理接口。
\end{itemize}

当前后端主路由由 \texttt{fake\_image\_detector/urls.py} 将 \texttt{core.urls} 挂载在 \texttt{/api/} 下，因此所有 HTTP 接口统一以 \texttt{/api/} 作为访问前缀；实时通知接口通过 WebSocket 路由 \texttt{/ws/notifications/} 暴露。

\subsection{内部接口分类}
系统内部接口主要包括：
\begin{itemize}[leftmargin=2em]
    \item \textbf{后端与异步任务系统接口}：由后端创建检测任务、投递 Celery 异步任务，并在任务完成后回写状态和生成通知；
    \item \textbf{后端与 AI 推理服务接口}：由图像检测模块、后续论文检测模块和 Review 检测模块调用对应推理服务；
    \item \textbf{后端与文件存储接口}：用于读取上传文件、提取图像、生成报告、保存中间产物和对外下载；
    \item \textbf{后端与消息通道接口}：用于向 WebSocket 消费者发送检测完成通知、审核任务通知和系统广播消息。
\end{itemize}

\subsection{统一接口约定}
为了保证接口风格一致，本系统采用如下约定：
\begin{itemize}[leftmargin=2em]
    \item HTTP 接口统一返回 JSON 结构，文件下载接口返回文件流；
    \item 认证接口返回访问令牌、刷新令牌或用户上下文，业务接口通过认证信息识别当前用户；
    \item 涉及文件上传的接口使用 \texttt{multipart/form-data}；
    \item 涉及列表展示的接口支持分页或分页式响应结构；
    \item 涉及任务状态的接口应明确返回 \texttt{pending}、\texttt{in\_progress}、\texttt{completed} 等状态字段；
    \item 错误响应应区分认证失败、权限不足、资源不存在、参数错误和系统异常。
\end{itemize}

\subsection{当前项目外部接口总览}
\begin{longtable}{|p{3.2cm}|p{2cm}|p{8.6cm}|}
\hline
接口类别 & 访问方式 & 说明 \\ \hline
用户认证接口 & HTTP & 提供注册、登录、登出、令牌刷新、用户资料查询与更新、头像更新和密码重置。 \\ \hline
组织管理接口 & HTTP & 提供组织申请、审批、组织详情、邀请码和组织权限更新。 \\ \hline
资源上传接口 & HTTP & 提供文件上传、文件详情查询、图片提取结果查询、文件标签更新和文件删除。 \\ \hline
检测任务接口 & HTTP & 提供图像检测提交、任务状态查询、任务结果列表、任务删除和报告下载。 \\ \hline
人工审核接口 & HTTP & 提供审核申请创建、审核任务列表、审核详情、审核提交、人工审核报告生成等能力。 \\ \hline
通知接口 & HTTP + WebSocket & 提供通知列表、未读数量、已读状态更新、广播通知和实时推送。 \\ \hline
管理端治理接口 & HTTP & 提供管理员登录、用户管理、文件管理、日志查询、仪表盘统计和审核审批。 \\ \hline
\end{longtable}

\subsection{用户认证与个人信息接口}
该类接口主要服务用户端与管理端登录态建立和基础资料维护。当前项目的关键接口如下：
\begingroup
\small
\setlength{\tabcolsep}{4pt}
\renewcommand{\arraystretch}{1.15}
\begin{longtable}{|L{6.2cm}|L{1.4cm}|L{5.6cm}|}
\hline
接口路径 & 方法 & 说明 \\ \hline
\endfirsthead
\hline
接口路径 & 方法 & 说明 \\ \hline
\endhead
\hline
\endfoot
\hline
\endlastfoot
\path{/api/register/} & POST & 用户注册。 \\ \hline
\path{/api/login/} & POST & 普通用户登录。 \\ \hline
\path{/api/logout/} & POST & 用户登出。 \\ \hline
\path{/api/token/refresh/} & POST & 刷新认证令牌。 \\ \hline
\path{/api/user/details/} & GET & 获取当前用户详情。 \\ \hline
\path{/api/user/update/} & PUT & 更新用户资料。 \\ \hline
\path{/api/user/avatar/} & PUT & 更新用户头像。 \\ \hline
\path{/api/password-reset/} & POST & 发送密码重置验证码。 \\ \hline
\path{/api/password-reset/confirm/} & POST & 提交验证码并完成密码重置。 \\ \hline
\path{/api/admin-login/} & POST & 管理员登录。 \\ \hline
\end{longtable}
\endgroup

从设计上看，这类接口构成所有业务接口的前置条件。后续在扩展统一资源与统一任务接口时，仍然需要依赖这一层提供用户身份、角色和组织上下文。

\subsection{组织与成员接口}
组织相关接口用于支撑组织入驻、组织审批、邀请码管理和组织权限维护。当前关键接口如下：
\begingroup
\small
\setlength{\tabcolsep}{4pt}
\renewcommand{\arraystretch}{1.15}
\begin{longtable}{|L{6.2cm}|L{1.4cm}|L{5.6cm}|}
\hline
接口路径 & 方法 & 说明 \\ \hline
\endfirsthead
\hline
接口路径 & 方法 & 说明 \\ \hline
\endhead
\hline
\endfoot
\hline
\endlastfoot
\path{/api/organization/create/} & POST & 提交组织申请。 \\ \hline
\path{/api/organizations/create-directly/} & POST & 管理端直接创建组织。 \\ \hline
\path{/api/organization/applications/get_pending/} & GET & 获取待审批的组织申请列表。 \\ \hline
\path{/api/organization/applications/<app_id>/} & GET & 查看组织申请详情。 \\ \hline
\path{/api/organization/<app_id>/approve/} & POST & 审批通过组织申请。 \\ \hline
\path{/api/organization/<app_id>/reject/} & POST & 拒绝组织申请。 \\ \hline
\path{/api/organizations/} & GET & 获取组织列表。 \\ \hline
\path{/api/organization/<org_id>/} & GET & 获取组织详情。 \\ \hline
\path{/api/organization/<org_id>/delete/} & DELETE & 删除组织。 \\ \hline
\path{/api/organization/<org_id>/permission/} & POST & 更新组织角色权限。 \\ \hline
\path{/api/organization/<org_id>/invitation_codes/} & GET & 获取邀请码列表。 \\ \hline
\path{/api/organization/usage/} & GET & 查询组织检测次数使用情况。 \\ \hline
\path{/api/organization/recharge-uses/} & POST & 为用户或组织补充检测次数。 \\ \hline
\end{longtable}
\endgroup

\subsection{资源上传与资源管理接口}
该类接口负责接收上传文件、读取文件详情、获取提取后的图像、维护文件标签和删除资源。当前实现仍以文件和图像资源为中心，后续论文与 Review 资源可在统一资源接口下复用同一接入逻辑。

\begingroup
\small
\setlength{\tabcolsep}{4pt}
\renewcommand{\arraystretch}{1.15}
\begin{longtable}{|L{6.2cm}|L{1.4cm}|L{5.6cm}|}
\hline
接口路径 & 方法 & 说明 \\ \hline
\endfirsthead
\hline
接口路径 & 方法 & 说明 \\ \hline
\endhead
\hline
\endfoot
\hline
\endlastfoot
\path{/api/upload/} & POST & 上传文件，支持图片、PDF、压缩包等。 \\ \hline
\path{/api/upload/<file_id>/} & GET & 获取单个上传文件详情。 \\ \hline
\path{/api/upload/<file_id>/extract_images/} & GET & 获取从文件中提取出的图像集合。 \\ \hline
\path{/api/upload/<file_id>/addTag/} & POST & 更新文件标签。 \\ \hline
\path{/api/upload/<file_id>/delete/} & DELETE & 删除上传文件。 \\ \hline
\path{/api/upload/get_all_file_images/<file_management_id>/} & GET & 获取某个文件下的全部图像。 \\ \hline
\path{/api/get_files/} & GET & 管理端获取文件列表。 \\ \hline
\path{/api/delete_image_upload/<image_id>/} & DELETE & 管理端删除单张图像资源。 \\ \hline
\end{longtable}
\endgroup

\subsection{检测任务与检测结果接口}
检测接口是当前项目接口体系的核心。现有实现已经形成图像检测任务的完整闭环，包括任务提交、状态查询、结果列表、详情查询和报告下载。后续论文检测与 Review 检测建议继续沿用统一任务接口设计。

\begingroup
\small
\setlength{\tabcolsep}{4pt}
\renewcommand{\arraystretch}{1.15}
\begin{longtable}{|L{6.2cm}|L{1.4cm}|L{5.6cm}|}
\hline
接口路径 & 方法 & 说明 \\ \hline
\endfirsthead
\hline
接口路径 & 方法 & 说明 \\ \hline
\endhead
\hline
\endfoot
\hline
\endlastfoot
\path{/api/detection/submit/} & POST & 提交图像检测任务。 \\ \hline
\path{/api/detection/<image_id>/} & GET & 获取单张图像的检测结果。 \\ \hline
\path{/api/detection-task/<task_id>/status/} & GET & 用户端查询检测任务状态。 \\ \hline
\path{/api/get_detection_task_status/<task_id>/} & GET & 管理端查询检测任务状态。 \\ \hline
\path{/api/tasks/<task_id>/results/} & GET & 获取任务全部结果。 \\ \hline
\path{/api/tasks/<task_id>/fake_results/} & GET & 获取任务中判定为可疑的结果。 \\ \hline
\path{/api/tasks/<task_id>/normal_results/} & GET & 获取任务中判定为正常的结果。 \\ \hline
\path{/api/results/<result_id>/} & GET & 获取单条检测结果详情。 \\ \hline
\path{/api/results_image/<image_id>/} & GET & 通过图像查询检测结果。 \\ \hline
\path{/api/tasks_image/<image_id>/getdr/} & GET & 通过图像获取检测结果映射。 \\ \hline
\path{/api/tasks/<task_id>/report/} & GET & 下载任务级检测报告。 \\ \hline
\path{/api/tasks_image/<image_id>/report/} & GET & 下载图像级检测报告。 \\ \hline
\path{/api/detection-task-delete/<task_id>/} & DELETE & 删除检测任务。 \\ \hline
\path{/api/user-tasks/} & GET & 获取当前用户检测任务列表。 \\ \hline
\path{/api/task-summary/} & GET & 获取当前用户任务摘要。 \\ \hline
\path{/api/get-task-summary/} & GET & 获取任务汇总信息。 \\ \hline
\end{longtable}
\endgroup

\subsection{人工审核接口}
人工审核接口构成自动检测之后的复核链路。该类接口覆盖审核申请、审核任务查看、审核结果查看、审核提交和人工审核报告生成。

\begingroup
\small
\setlength{\tabcolsep}{4pt}
\renewcommand{\arraystretch}{1.15}
\begin{longtable}{|L{6.2cm}|L{1.4cm}|L{5.6cm}|}
\hline
接口路径 & 方法 & 说明 \\ \hline
\endfirsthead
\hline
接口路径 & 方法 & 说明 \\ \hline
\endhead
\hline
\endfoot
\hline
\endlastfoot
\path{/api/create_review_task_with_admin_check/} & POST & 创建人工审核申请，等待管理员审批。 \\ \hline
\path{/api/get_request_completion_status/<task_id>/} & GET & 查询某检测任务对应审核申请的完成情况。 \\ \hline
\path{/api/get_request_detail/<reviewRequest_id>/} & GET & 获取审核申请详情。 \\ \hline
\path{/api/get_publisher_review_tasks/} & GET & 发布者查看本人发起的审核任务。 \\ \hline
\path{/api/get_reviewer_tasks/} & GET & 审核人员获取待处理审核任务列表。 \\ \hline
\path{/api/get-reviewer-request-detail/<reviewRequest_id>/} & GET & 审核人员查看审核申请详情。 \\ \hline
\path{/api/get_review_detail/<manual_review_id>/} & GET & 获取单个审核任务详情。 \\ \hline
\path{/api/post_review/<manual_review_id>/} & POST & 提交人工审核结果。 \\ \hline
\path{/api/get_img_review_all/} & GET & 获取审核任务下所有图像审核结果。 \\ \hline
\path{/api/get_image_review/} & GET & 获取指定图像的人工审核结果。 \\ \hline
\path{/api/manual-review/<review_id>/report/} & GET & 生成或下载人工审核报告。 \\ \hline
\path{/api/manual-review/<review_request_id>/} & GET & 通过审核申请获取关联人工审核记录。 \\ \hline
\path{/api/publisher-dectectiontask-access/} & GET & 检查发布者是否有权访问某检测任务。 \\ \hline
\path{/api/reviewer-manualreview-access/} & GET & 检查审核人员是否有权访问某人工审核记录。 \\ \hline
\path{/api/publishers/<publisher_id>/reviewers/} & GET & 获取发布者可关联审核人员。 \\ \hline
\path{/api/get_all_reviewers/} & GET & 获取组织内审核人员列表。 \\ \hline
\end{longtable}
\endgroup

\subsection{通知与实时消息接口}
通知接口分为持久化通知接口和实时推送接口两类。持久化通知负责通知列表和已读状态维护，实时推送接口负责检测完成等事件的在线通知。

\begingroup
\small
\setlength{\tabcolsep}{4pt}
\renewcommand{\arraystretch}{1.15}
\begin{longtable}{|L{6.2cm}|L{1.4cm}|L{5.6cm}|}
\hline
接口路径 & 方法 & 说明 \\ \hline
\endfirsthead
\hline
接口路径 & 方法 & 说明 \\ \hline
\endhead
\hline
\endfoot
\hline
\endlastfoot
\path{/api/notification/get/} & GET & 获取当前用户通知列表。 \\ \hline
\path{/api/notification/notify/} & GET & 获取当前用户未读通知数量。 \\ \hline
\path{/api/notification/set_as_read/} & POST & 将当前用户全部通知设为已读。 \\ \hline
\path{/api/notification/set_as_read/<notification_id>/} & POST & 将单条通知设为已读。 \\ \hline
\path{/api/notification/broadcast/} & POST & 管理端发送广播通知。 \\ \hline
\path{/ws/notifications/} & WebSocket & 实时推送通知与任务状态更新。 \\ \hline
\end{longtable}
\endgroup

\subsection{管理端治理与统计接口}
管理端接口覆盖管理员资料、仪表盘统计、用户管理、文件管理、审核审批和日志查询等能力。该类接口支撑当前项目的治理能力，也是后续扩展论文检测日志、Review 检测日志和模型配置管理的入口。

\begingroup
\small
\setlength{\tabcolsep}{4pt}
\renewcommand{\arraystretch}{1.15}
\begin{longtable}{|L{6.2cm}|L{1.4cm}|L{5.6cm}|}
\hline
接口路径 & 方法 & 说明 \\ \hline
\endfirsthead
\hline
接口路径 & 方法 & 说明 \\ \hline
\endhead
\hline
\endfoot
\hline
\endlastfoot
\path{/api/admin/details/} & GET & 获取当前管理员详情。 \\ \hline
\path{/api/admin/details/<user_id>} & GET & 获取指定用户详情。 \\ \hline
\path{/api/admin_dashboard/} & GET & 获取管理端仪表盘汇总数据。 \\ \hline
\path{/api/dashboard/img_tag/} & GET & 图像标签统计。 \\ \hline
\path{/api/dashboard/top_publishers/} & GET & 发布者排行统计。 \\ \hline
\path{/api/dashboard/top_organizations/} & GET & 组织排行统计。 \\ \hline
\path{/api/dashboard/daily_active_users/} & GET & 日活用户统计。 \\ \hline
\path{/api/dashboard/daily_active_organizations/} & GET & 日活组织统计。 \\ \hline
\path{/api/dashboard/daily_task_count/} & GET & 每日检测任务统计。 \\ \hline
\path{/api/dashboard/daily_review_request_count/} & GET & 每日审核申请统计。 \\ \hline
\path{/api/dashboard/daily_completed_manual_review_count/} & GET & 每日完成审核统计。 \\ \hline
\path{/api/dashboard/get_sub_method_distribution_by_tag/} & GET & 子检测方法分布统计。 \\ \hline
\path{/api/get_users/} & GET & 获取用户列表。 \\ \hline
\path{/api/create_user/} & POST & 管理端创建用户。 \\ \hline
\path{/api/update_user/<user_id>/} & PUT & 更新用户。 \\ \hline
\path{/api/delete_user/<user_id>/} & DELETE & 删除用户。 \\ \hline
\path{/api/create-admin/} & POST & 创建管理员。 \\ \hline
\path{/api/user_permission/<user_id>/} & POST & 更新用户权限。 \\ \hline
\path{/api/manage-associations/} & POST & 建立发布者与审核者关联。 \\ \hline
\path{/api/user_action_log/} & GET & 查询用户操作日志。 \\ \hline
\path{/api/user_action_log/<log_id>/} & DELETE & 删除日志。 \\ \hline
\path{/api/user_action_log/download/} & GET & 导出日志。 \\ \hline
\path{/api/get_reviewRequest/all/} & GET & 获取全部审核申请。 \\ \hline
\path{/api/get_reviewRequest/<reviewRequest_id>/} & GET & 获取审核申请详情。 \\ \hline
\path{/api/handle_reviewRequest/<reviewRequest_id>/} & POST & 管理员处理审核申请。 \\ \hline
\path{/api/get_review_request_detail/<manual_review_id>/} & GET & 查看人工审核执行详情。 \\ \hline
\path{/api/post_report/<post_id>/} & POST & 处理帖子举报。 \\ \hline
\path{/api//review-requests/<review_request_id>/delete/} & DELETE & 删除审核申请；当前代码路径定义带前导斜杠，建议后续修正。 \\ \hline
\end{longtable}
\endgroup

\subsection{内部接口设计}
\subsubsection{后端与异步任务系统接口}
当前系统通过检测提交接口创建任务后，调用 Celery 执行异步检测。该类内部接口主要承担以下职责：
\begin{itemize}[leftmargin=2em]
    \item 接收检测任务 ID、图像 ID 集合、检测参数和用户上下文；
    \item 在异步执行中更新任务状态、检测结果、报告生成状态和通知状态；
    \item 在任务完成后通过消息通道向 WebSocket 消费者推送结果。
\end{itemize}

\subsubsection{后端与 AI 推理服务接口}
当前项目的图像检测能力通过后端调用图像推理逻辑实现。后续论文检测和 Review 检测建议采用与图像检测一致的适配方式：
\begin{itemize}[leftmargin=2em]
    \item 后端传入资源路径、文本内容或结构化输入；
    \item 推理层返回风险评分、分类结论、可视化结果或结构化分析结果；
    \item 后端统一完成结果落库、状态更新、报告生成和通知触发。
\end{itemize}

\subsubsection{后端与文件存储接口}
文件存储接口不直接对前端暴露，而由业务服务封装。主要职责包括：
\begin{itemize}[leftmargin=2em]
    \item 保存上传文件、提取图像、检测中间产物和报告文件；
    \item 提供受控读取能力，供结果查询和报告下载接口调用；
    \item 在删除任务或删除资源时执行对应文件清理。
\end{itemize}

\subsubsection{后端与消息通道接口}
后端通过 Django Channels 将状态变化推送到 WebSocket 通道。内部消息接口应保证：
\begin{itemize}[leftmargin=2em]
    \item 推送事件与数据库通知记录保持一致；
    \item 用户上线和离线不影响通知持久化；
    \item 前端可基于同一消息体刷新未读数、任务状态或详情页数据。
\end{itemize}

\subsection{接口安全设计}
接口安全主要包括以下方面：
\begin{itemize}[leftmargin=2em]
    \item 需要登录的接口必须通过认证令牌识别当前用户；
    \item 用户只能访问本人资源、本人检测任务、本人通知和权限范围内的审核记录；
    \item 管理端接口必须校验管理员身份和管理权限；
    \item 文件下载接口必须校验资源归属和访问权限，避免直接暴露物理路径；
    \item 审核、权限调整、组织审批、日志导出等高风险接口必须记录审计日志。
\end{itemize}

\subsection{接口演进方向}
从当前项目代码看，接口仍然明显偏向图像检测场景。后续接口演进建议如下：
\begin{itemize}[leftmargin=2em]
    \item 引入统一资源接口，将图像、论文和 Review 统一纳入资源上传和资源查询入口；
    \item 引入统一检测任务接口，使图像检测、论文检测和 Review 检测共用任务提交和状态查询协议；
    \item 保留差异化结果接口，分别承载图像可疑区域、论文段落概率和 Review 模板化倾向等差异结果；
    \item 将论文检测日志、Review 检测日志和模型配置接口纳入管理端统一治理接口空间；
    \item 对现有路径命名、重复路由和个别异常路径进行收敛，提升整体接口规范性。
\end{itemize}

\subsection{用户端页面框架与导航路由设计}\label{sec:user-client-routing}
用户端基于 Vue~3 与 Vuetify，采用 \texttt{vue-router} 文件式路由（\texttt{src/pages} 目录映射 URL）。迭代阶段需在保留现有图像上传、检测历史、人工审核与审阅员待办等能力的前提下，补齐\textbf{全篇论文 AIGC 检测}、\textbf{同行评审 Review 文本检测}与\textbf{造假检测记录统一管理}的页面骨架及导航入口，使三类检测对象在交互结构上对齐、在路由上可区分、在记录侧可汇总。

\subsubsection{全局壳层与导航原则}
\begin{itemize}[leftmargin=2em]
    \item \textbf{布局壳层}：继续使用应用根布局中的顶栏、侧栏（桌面端）与底部导航（移动端）承载一级入口；业务页面仅负责 \texttt{router-view} 内主内容区，不在各业务页重复实现顶栏与认证态判断。
    \item \textbf{鉴权与角色}：沿用全局路由守卫，未登录用户统一跳转登录页；发布者（\texttt{publisher}）可见检测与记录相关入口，审核者（\texttt{reviewer}）可见人工审阅待办入口，与需求规格中用户展示与权限划分一致。
    \item \textbf{路由命名区分}：现有路径 \path{/review} 表示审核参与者「我的审阅任务」，与「同行评审 Review 文本鉴伪检测」业务不同；迭代中 Review 文本检测应使用独立路径前缀（如 \path{/detect/review}），避免语义冲突与后期维护混淆。
    \item \textbf{异步任务一致体验}：论文与 Review 检测均为长时任务，页面框架须预留任务 ID、状态轮询或 WebSocket 订阅、错误与重试区域，与图像检测任务交互模式保持一致，降低用户学习成本。
\end{itemize}

\subsubsection{全篇论文 AIGC 检测页面框架}
\textbf{目标：}支撑用户上传论文文件（如 PDF、DOCX）、填写任务名称与检测选项、提交后查看排队/进行中/完成状态，并在完成后跳转结果摘要或统一记录详情。

\textbf{建议路由：}
\begin{itemize}[leftmargin=2em]
    \item \path{/detect/paper}：论文检测主页面（上传区 + 参数区 + 提交区）；
    \item \path{/detect/paper/task/:taskId}：单任务进度与结果入口（可与记录详情复用同一详情布局，仅默认来源为「从论文检测进入」）。
\end{itemize}

\textbf{页面分区（框架级）：}
\begin{enumerate}[leftmargin=2em]
    \item \textbf{上传与格式说明区}：文件选择、大小与格式提示、组织配额或次数提示（与后端 \path{/api/organization/usage/} 等接口衔接，具体字段以实现为准）；
    \item \textbf{任务元数据区}：任务名称、是否启用事实性鉴伪等开关（与论文检测模块参数对齐）；
    \item \textbf{状态与操作区}：提交按钮、进行中进度、失败重试与「前往检测记录」链接；
    \item \textbf{预留结果摘要区}：任务完成后展示全文级风险、段落级入口或跳转详情的按钮（详细可视化在详细设计阶段展开）。
\end{enumerate}

\begin{figure}[H]
    \centering
    \resizebox{0.98\textwidth}{!}{\begin{tikzpicture}[
    font=\scriptsize,
    box/.style={draw, rounded corners=2pt, line width=0.6pt},
    sec/.style={draw, rounded corners=2pt, line width=0.6pt, fill=gray!6},
    lbl/.style={anchor=west, align=left},
]
% Page frame
\draw[box] (0,0) rectangle (16,9.2);
\node[lbl] at (0.4,8.85) {\textbf{全篇论文 AIGC 检测页（/detect/paper）}};

% Header / breadcrumbs
\draw[sec] (0.6,8.15) rectangle (15.4,8.55);
\node[lbl, text width=14.2cm] at (0.9,8.35) {导航/标题区：论文检测 \quad|\quad 任务状态 \quad|\quad 帮助说明};

% Left: upload + meta
\draw[sec] (0.6,4.9) rectangle (10.2,7.95);
\node[lbl] at (0.9,7.65) {\textbf{A. 上传与格式说明}};
\node[lbl, text width=9.0cm] at (0.9,7.25) {文件选择（PDF/DOCX/TXT） \quad|\quad 大小/格式校验提示};
\node[lbl, text width=9.0cm] at (0.9,6.85) {组织/个人配额提示 \quad|\quad 隐私与合规提示};
\draw[box] (0.9,5.35) rectangle (9.9,6.55);
\node at (5.4,5.95) {拖拽上传区 / 文件列表（占位）};

\draw[sec] (0.6,1.45) rectangle (10.2,4.6);
\node[lbl] at (0.9,4.30) {\textbf{B. 任务元数据与检测选项}};
\node[lbl] at (0.9,3.90) {任务名称 \quad| \quad 学科标签（可选）};
\node[lbl, text width=9.0cm] at (0.9,3.50) {开关：事实性鉴伪（可选） \quad|\quad 模型/阈值（占位）};
\draw[box] (0.9,1.85) rectangle (9.9,3.10);
\node at (5.4,2.45) {参数表单区域（占位）};

% Right: submit + status + summary
\draw[sec] (10.6,4.9) rectangle (15.4,7.95);
\node[lbl] at (10.9,7.65) {\textbf{C. 提交与状态}};
\draw[box] (10.9,7.10) rectangle (15.1,7.42);
\node at (13.0,7.26) {提交检测 / 取消};
\draw[box] (10.9,5.35) rectangle (15.1,6.95);
\node[align=center, text width=4.0cm] at (13.0,6.15) {状态卡片：排队/进行中/完成/失败\\进度条 + 重试 + 前往记录};

\draw[sec] (10.6,1.45) rectangle (15.4,4.6);
\node[lbl] at (10.9,4.30) {\textbf{D. 结果摘要（完成后显示）}};
\draw[box] (10.9,1.85) rectangle (15.1,4.00);
\node[align=center, text width=4.0cm] at (13.0,2.95) {全文风险等级 + AI贡献占比（占位）\\高风险段落入口/报告下载入口};

% Footer actions
\draw[sec] (0.6,0.45) rectangle (15.4,1.10);
\node[lbl, text width=14.2cm] at (0.9,0.78) {页脚辅助：FAQ \quad|\quad 反馈入口 \quad|\quad 任务提示（可选）};
\end{tikzpicture}

}
    \caption{用户端全篇论文 AIGC 检测页线框图（最小框架）}
\end{figure}

\subsubsection{同行评审 Review 检测页面框架}
\textbf{目标：}支撑用户粘贴 Review 文本或上传含评审意见的文件，发起针对「虚假评审/模板化/AI 生成倾向」的检测，并复用统一任务状态与通知机制。

\textbf{建议路由：}
\begin{itemize}[leftmargin=2em]
    \item \path{/detect/review}：Review 检测主页面（文本编辑区或文件上传区二选一或并列）；
    \item \path{/detect/review/task/:taskId}：单任务进度与结果入口（同样可与记录详情复用）。
\end{itemize}

\textbf{页面分区（框架级）：}
\begin{enumerate}[leftmargin=2em]
    \item \textbf{输入模式切换}：纯文本输入与文件上传切换，便于短评与长文档两种场景；
    \item \textbf{检测参数区}：语言、阈值或策略占位（与评审检测模块后续接口字段对齐）；
    \item \textbf{证据展示占位}：为可疑片段高亮、模板化提示预留主内容槽位，框架阶段可用占位组件标出数据流挂载点；
    \item \textbf{合规提示区}：提示用户仅检测有权处理的评审文本，符合学术数据使用边界。
\end{enumerate}

\begin{figure}[H]
    \centering
    \resizebox{0.98\textwidth}{!}{\begin{tikzpicture}[
    font=\scriptsize,
    box/.style={draw, rounded corners=2pt, line width=0.6pt},
    sec/.style={draw, rounded corners=2pt, line width=0.6pt, fill=gray!6},
    chip/.style={draw, rounded corners=9pt, line width=0.6pt, fill=blue!6},
    lbl/.style={anchor=west, align=left},
]
\draw[box] (0,0) rectangle (16,9.2);
\node[lbl] at (0.4,8.85) {\textbf{Review 文本检测页（/detect/review）}};

% Header
\draw[sec] (0.6,8.15) rectangle (15.4,8.55);
\node[lbl, text width=14.2cm] at (0.9,8.35) {导航/标题区：Review 检测 \quad|\quad 合规提示};

% Mode switch chips
\node[chip, minimum width=3.2cm, minimum height=0.55cm] (m1) at (3.0,7.55) {文本输入};
\node[chip, minimum width=3.2cm, minimum height=0.55cm, fill=gray!10] (m2) at (6.6,7.55) {文件上传};
\node[lbl] at (0.9,7.55) {\textbf{输入模式：}};

% Left main input
\draw[sec] (0.6,3.05) rectangle (10.6,7.20);
\node[lbl] at (0.9,6.90) {\textbf{A. Review 输入区}};
\draw[box] (0.9,3.55) rectangle (10.3,6.55);
\node[align=center, text width=8.9cm] at (5.6,5.05) {多行文本编辑器/文件解析预览（占位）\\支持粘贴、字数统计、清洗提示};
\draw[box] (0.9,3.25) rectangle (10.3,3.45);
\node[lbl] at (1.1,3.35) {提示：仅检测有权处理的评审文本};

% Right: parameters
\draw[sec] (10.9,5.10) rectangle (15.4,7.20);
\node[lbl] at (11.2,6.90) {\textbf{B. 检测参数（占位）}};
\node[lbl] at (11.2,6.50) {语言 / 规则版本 / 阈值};
\draw[box] (11.2,5.35) rectangle (15.1,6.25);
\node at (13.15,5.80) {参数表单};

% Right: submit + status
\draw[sec] (10.9,3.05) rectangle (15.4,4.95);
\node[lbl] at (11.2,4.65) {\textbf{C. 提交与状态}};
\draw[box] (11.2,4.20) rectangle (15.1,4.52);
\node at (13.15,4.36) {提交检测 / 清空};
\draw[box] (11.2,3.35) rectangle (15.1,4.05);
\node[align=center, text width=3.8cm] at (13.15,3.70) {状态：排队/进行中/完成/失败\\重试 + 前往记录};

% Evidence / highlight area
\draw[sec] (0.6,0.45) rectangle (15.4,2.75);
\node[lbl] at (0.9,2.45) {\textbf{D. 证据与解释展示（完成后显示）}};
\draw[box] (0.9,0.85) rectangle (15.1,2.20);
\node[align=center, text width=14.0cm] at (8.0,1.52) {风险评分 + 结论摘要（占位）\quad|\quad 可疑片段高亮/模板化提示（占位）};
\end{tikzpicture}

}
    \caption{用户端 Review 文本检测页线框图（最小框架）}
\end{figure}

\subsubsection{造假检测记录管理页面框架与路由}
\textbf{目标：}对应需求中的「造假检测记录管理」：用户按时间查看本人发起的检测任务，覆盖\textbf{学术图像、全篇论文、Review}等类型，支持筛选、检索、进入详情与报告下载（与 FR-YHZS-0002 一致）。

\textbf{路由策略：}
\begin{itemize}[leftmargin=2em]
    \item \textbf{演进方案 A（兼容现有）}：保留现有 \path{/history} 作为统一「检测记录」入口，在页面内以 Tab 或分段列表区分任务类型（图像 / 论文 / Review），列表项展示任务类型标签、状态、时间与报告入口；
    \item \textbf{演进方案 B（路径语义化）}：新增 \path{/records} 作为统一记录中心，\path{/history} 通过路由重定向至 \path{/records}，减少后续文档与代码中的「历史」与「全类型记录」双轨表述。
\end{itemize}

\textbf{建议子路由（记录详情，可与各检测子页复用）：}
\begin{itemize}[leftmargin=2em]
    \item \path{/records/:taskId} 或 \path{/history/:taskId}：单条检测记录详情（状态时间线、类型专属结果槽位、报告下载）；
    \item 查询参数占位：在记录列表路径上附加查询参数 \texttt{type}（取值如 \texttt{paper}、\texttt{image}、\texttt{review}）支持类型筛选，便于导航栏「记录」直达某类列表。
\end{itemize}

\textbf{页面分区（框架级）：}
\begin{enumerate}[leftmargin=2em]
    \item \textbf{筛选与检索条}：时间范围、内容类型、关键字（任务名等）；
    \item \textbf{列表区}：分页或虚拟滚动列表，列含类型、名称、状态、时间、操作（详情、下载、删除若开放）；
    \item \textbf{空状态与加载状态}：无数据与接口失败时的统一占位；
    \item \textbf{详情抽屉或独立详情页}：与 \path{/api/user-tasks/}、\path{/api/tasks/<task_id>/results/} 等统一任务接口配合，详情内按 \texttt{task\_type} 或等价字段切换结果子视图。
\end{enumerate}

\begin{figure}[H]
    \centering
    \resizebox{0.98\textwidth}{!}{\begin{tikzpicture}[
    font=\scriptsize,
    box/.style={draw, rounded corners=2pt, line width=0.6pt},
    sec/.style={draw, rounded corners=2pt, line width=0.6pt, fill=gray!6},
    pill/.style={draw, rounded corners=10pt, line width=0.6pt, fill=gray!8},
    lbl/.style={anchor=west, align=left},
]
\draw[box] (0,0) rectangle (16,9.2);
\node[lbl] at (0.4,8.85) {\textbf{检测记录列表页（/history 或 /records）}};

% Header
\draw[sec] (0.6,8.15) rectangle (15.4,8.55);
\node[lbl] at (0.9,8.35) {标题：检测记录 \quad|\quad 统计摘要（可选）};

% Filters row
\draw[sec] (0.6,7.10) rectangle (15.4,7.95);
\node[lbl] at (0.9,7.70) {\textbf{筛选与检索：}};
\draw[pill] (3.0,7.25) rectangle (5.2,7.80);
\node at (4.1,7.52) {类型：全部};
\draw[pill] (5.5,7.25) rectangle (8.3,7.80);
\node at (6.9,7.52) {时间：近半年};
\draw[pill] (8.6,7.25) rectangle (12.0,7.80);
\node at (10.3,7.52) {关键字：任务名};
\draw[pill] (12.3,7.25) rectangle (13.7,7.80);
\node at (13.0,7.52) {查询};
\draw[pill] (14.0,7.25) rectangle (15.2,7.80);
\node at (14.6,7.52) {重置};

% Table header
\draw[sec] (0.6,6.30) rectangle (15.4,6.90);
\node[lbl, text width=14.2cm] at (0.9,6.58) {列：类型 \quad|\quad 任务名称 \quad|\quad 状态 \quad|\quad 时间 \quad|\quad 操作（详情/报告/删除）};

% List area with rows
\draw[box] (0.6,1.35) rectangle (15.4,6.20);
\foreach \y/\t in {5.75/1,5.05/2,4.35/3,3.65/4,2.95/5,2.25/6,1.55/7} {
  \draw (0.6,\y) -- (15.4,\y);
}
\node[lbl, text width=14.2cm] at (0.9,5.95) {[图像] \quad 任务A \quad|\quad completed \quad|\quad 2026-04-19 \quad|\quad 详情/下载};
\node[lbl, text width=14.2cm] at (0.9,5.25) {[\textbf{论文}] \quad 任务B \quad|\quad in\_progress \quad|\quad 2026-04-18 \quad|\quad 查看进度};
\node[lbl, text width=14.2cm] at (0.9,4.55) {[\textbf{Review}] \quad 任务C \quad|\quad failed \quad|\quad 2026-04-17 \quad|\quad 重试};
\node[lbl] at (0.9,3.85) {…（更多记录）};

% Pagination
\draw[sec] (0.6,0.45) rectangle (15.4,1.10);
\node[lbl, text width=14.2cm] at (0.9,0.78) {分页：上一页 \quad 1 / N \quad 下一页 \quad|\quad 每页条数};
\end{tikzpicture}

}
    \caption{用户端检测记录列表页线框图（最小框架）}
\end{figure}

\subsubsection{检测记录详情页框架（复用统一任务详情）}
检测记录详情页用于承载三类检测对象的统一任务状态展示，并在同一页面内通过“类型专属结果槽位”挂载图像可疑区域、论文段落级结果或 Review 可疑片段等差异化内容。该页面应提供报告下载与可选的删除、人工审核入口，使用户能从“记录中心”闭环进入报告与复核链路。

\begin{figure}[H]
    \centering
    \resizebox{0.98\textwidth}{!}{\begin{tikzpicture}[
    font=\scriptsize,
    box/.style={draw, rounded corners=2pt, line width=0.6pt},
    sec/.style={draw, rounded corners=2pt, line width=0.6pt, fill=gray!6},
    lbl/.style={anchor=west, align=left},
    tag/.style={draw, rounded corners=10pt, line width=0.6pt, fill=blue!6, minimum height=0.45cm},
]
\draw[box] (0,0) rectangle (16,9.2);
\node[lbl] at (0.4,8.85) {\textbf{检测记录详情页（/records/:taskId 或 /history/:taskId）}};

% Title + meta
\draw[sec] (0.6,8.10) rectangle (15.4,8.75);
\node[lbl, text width=12.6cm] at (0.9,8.40) {任务名称 + ID \quad|\quad 创建时间 \quad|\quad 发起人 \quad|\quad 组织（可选）};
\node[tag, minimum width=1.6cm] at (14.3,8.42) {论文/图像/Review};

% Timeline/status
\draw[sec] (0.6,6.55) rectangle (15.4,8.00);
\node[lbl] at (0.9,7.70) {\textbf{A. 状态时间线（统一任务框架）}};
\draw[box] (0.9,6.85) rectangle (15.1,7.45);
\node[align=center, text width=13.8cm] at (8.0,7.15) {pending $\rightarrow$ in\_progress $\rightarrow$ completed/failed \quad|\quad 进度/耗时/错误信息（占位）};
\node[lbl, text width=14.2cm] at (0.9,6.70) {操作：重试（失败时）\quad|\quad 返回列表 \quad|\quad 前往报告};

% Left: type-specific main content
\draw[sec] (0.6,1.45) rectangle (10.6,6.35);
\node[lbl] at (0.9,6.05) {\textbf{B. 类型专属结果区（槽位）}};
\draw[box] (0.9,2.10) rectangle (10.3,5.75);
\node[align=center, text width=9.1cm] at (5.6,3.95) {图像：可疑区域/子方法摘要（占位）\\论文：全文结论 + 段落列表入口（占位）\\Review：风险证据片段高亮（占位）};
\draw[box] (0.9,1.75) rectangle (10.3,2.00);
\node[lbl] at (1.1,1.87) {通用字段：风险等级、摘要说明、模型版本（占位）};

% Right: actions + attachments
\draw[sec] (10.9,1.45) rectangle (15.4,6.35);
\node[lbl] at (11.2,6.05) {\textbf{C. 操作与资源}};
\draw[box] (11.2,5.45) rectangle (15.1,5.75);
\node at (13.15,5.60) {下载报告};
\draw[box] (11.2,5.05) rectangle (15.1,5.35);
\node at (13.15,5.20) {申请人工审核（可选）};
\draw[box] (11.2,4.65) rectangle (15.1,4.95);
\node at (13.15,4.80) {删除记录（可选）};

\draw[box] (11.2,2.35) rectangle (15.1,4.35);
\node[align=center, text width=3.8cm] at (13.15,3.35) {关联资源：原文件/文本摘要（占位）\\附件：提取图像/片段（占位）};
\draw[box] (11.2,1.75) rectangle (15.1,2.15);
\node[align=center] at (13.15,1.95) {通知与日志入口（可选）};

% Footer
\draw[sec] (0.6,0.45) rectangle (15.4,1.10);
\node[lbl, text width=14.2cm] at (0.9,0.78) {页脚提示：结果仅供参考 \quad|\quad 问题反馈};
\end{tikzpicture}

}
    \caption{用户端检测记录详情页线框图（最小框架）}
\end{figure}

\subsubsection{主导航与路由配置一览}
下表给出迭代目标下用户端一级导航与路径的\textbf{概要设计约定}（实现时可按团队规范微调路径字符串，但应保持「审阅任务」与「Review 文本检测」路径分离）。

\begingroup
\small
\setlength{\tabcolsep}{4pt}
\renewcommand{\arraystretch}{1.2}
\begin{longtable}{|p{3.4cm}|p{4.2cm}|p{6.4cm}|}
\hline
导航名称（建议） & 路径 & 说明 \\ \hline
\endfirsthead
\hline
导航名称（建议） & 路径 & 说明 \\ \hline
\endhead
\hline
\endfoot
\hline
\endlastfoot
主页 & \path{/} & 欢迎与功能导览，可补充三类检测入口卡片。 \\ \hline
图像检测（现有） & \path{/upload}、\path{/step/:id} 等 & 保持现有上传与步骤流；导航标题可逐步改为「图像检测」以与论文、Review 并列。 \\ \hline
论文 AIGC 检测 & \path{/detect/paper}（及任务子路径） & 发布者可见；对接统一检测任务创建接口的论文类型。 \\ \hline
Review 文本检测 & \path{/detect/review}（及任务子路径） & 发布者可见；与 \path{/review}（审阅员待办）严格区分。 \\ \hline
检测记录 & \path{/history} 或 \path{/records} & 统一造假检测记录管理；含筛选与详情框架。 \\ \hline
人工审核（现有） & \path{/annual} & 保持现网路径；迭代中在文案与详情中扩展对论文、Review 审核的展示占位。 \\ \hline
我的审阅（现有） & \path{/review} & 审核者任务列表，非 Review 文本自动检测。 \\ \hline
个人主页 & \path{/profile} & 用户资料与设置。 \\ \hline
登录 & \path{/login} & 未登录入口。 \\ \hline
\end{longtable}
\endgroup

\subsubsection{与后端接口的衔接说明}
论文检测、Review 检测与统一记录在接口层应逐步收敛到「统一任务提交 + 统一状态查询 + 差异化结果查询」协议（见本章前文演进方向）；用户端路由仅负责\textbf{按任务类型与任务 ID 组织视图}，不在前端硬编码具体模型或算法版本。检测记录页应优先消费带任务类型与分页的任务列表接口，避免为图像、论文、Review 各写一套列表页而增加导航与权限维护成本。

\subsubsection{本小节与实现关系}
当前仓库用户端已实现图像上传、检测历史、人工审核与审阅任务等页面及 \path{/history}、\path{/upload}、\path{/review} 等路由；本小节给出论文与 Review 检测的\textbf{页面框架与路由配置概要}，作为详细设计与编码迭代的依据。落地时需在 \texttt{App.vue}（或等价布局）中为发布者增加「论文检测」「Review 检测」导航项，并实现对应 \texttt{pages} 下路由文件，使导航、路由与统一记录中心三者一致。

\subsection{本章小结}
本章从概要设计角度给出了系统接口体系的总体设计，并补充了\textbf{用户端}在论文 AIGC 检测、Review 文本检测与造假检测记录管理方面的页面框架与导航路由约定。当前项目已经形成以用户认证、资源上传、图像检测、人工审核、通知和管理端治理为主体的接口闭环；后续论文检测与 Review 检测能力可在统一资源、统一任务和差异化结果接口框架下继续扩展，并与统一检测记录路由协同演进。接口设计的核心目标不是增加更多零散接口，而是在现有工程基础上形成可持续演进的统一接口体系。
